\documentclass{article}
\usepackage[utf8]{inputenc}
\usepackage{graphicx}
\usepackage{amsmath}
\usepackage{amssymb}
\usepackage{hyperref}
\usepackage{geometry}

\geometry{margin=1in}

\title{Minimal Seed Structure: Emergent Intelligence from Local Connection and Minimal Modification}
\author{Chen Leiyang}
\date{\today}

\begin{document}

\maketitle

\begin{abstract}
We discover and validate a minimal seed structure for emergent intelligence: LocalConnection(k=4) + Prediction Loss + $\lambda \approx 0$. Through systematic exploration across multiple dimensions (activation functions, initialization, optimizer, topology, loss types, learning rates), we find that this minimal seed can generate all observed capabilities (causal emergence, identity tracking, compositional generalization, relational reasoning) through just 1-2 minimal modifications. We further discover the Constraint Interaction Law: capability emergence threshold = $\lambda \times Density / nonlinearity\_exponent$. We also find physics-dependent scaling: linear force systems prefer MLP while nonlinear force systems prefer GNN at larger N. Our findings suggest that intelligence may emerge from extremely simple structures through minimal modifications, challenging the "bigger is better" paradigm in modern AI.
\end{abstract}

\section{Introduction}

Modern AI systems rely heavily on massive parameters and compute. However, biological intelligence often emerges from simple structures through minimal modifications. This work explores the fundamental question:

\textbf{What is the minimal structure that can generate intelligence?}

Our contributions:
\begin{enumerate}
\item Discover Minimal Seed Structure: LocalConnection(k=4) + Prediction Loss + $\lambda \approx 0$
\item Systematically validate capability emergence through minimal modifications (L1-L9 framework)
\item Discover Constraint Interaction Law: $\lambda \times Density / nonlinearity\_exponent$ determines capability emergence
\item Find physics-dependent scaling: linear forces prefer MLP, nonlinear forces prefer GNN
\item Explore optimal configurations across dimensions (activation, topology, optimizer, etc.)
\end{enumerate}

\section{Related Work}

\subsection{Causal Representation Learning}
Prior work on causal representation learning focuses on discovering underlying factors from observations. Our work complements this by identifying the minimal structural requirements for causal emergence.

\subsection{Object-Centric Representations}
Slot Attention, SPACE, and similar methods aim to discover objects from images. We find that object identity tracking requires specific structural constraints (temporal binding).

\subsection{Emergence in Minimal Models}
Similar to how the minimal genome contains the essential information for life, or Ising model captures essential physics, we propose that Minimal Seed Structure captures essential requirements for intelligence.

\section{Minimal Seed Structure}

\subsection{Definition}
Our discovered minimal seed structure is:

\begin{equation}
\text{Seed} = \text{LocalConnection}(k=4) + \mathcal{L}_{prediction} + \lambda \approx 0
\end{equation}

Where:
\begin{itemize}
\item \textbf{LocalConnection(k=4)}: Each node connects to exactly 4 neighbors
\item \textbf{Prediction Loss}: Next-step position/velocity prediction
\item \textbf{$\lambda \approx 0$}: Minimal compression (free growth)
\end{itemize}

\subsection{Why k=4?}
Through systematic exploration of topology variants:
\begin{center}
\begin{tabular}{l|c}
Topology & MSE \\
\hline
Grid & 0.000086 \\
Ring & 0.000103 \\
SmallWorld p=0.1 & 0.000103 \\
MLP (no topo) & 0.000107 \\
SmallWorld p=0.3 & 0.000257 \\
Fully Connected & 0.003261 \\
\end{tabular}
\end{center}

Grid topology (k=4 neighbors in 2D) emerges as optimal.

\section{L1-L9 Capability Framework}

We discover that capabilities emerge through minimal modifications from the seed:

\begin{center}
\begin{tabular}{c|l|l|c}
Level & Capability & Minimal Modification & Trigger \\
\hline
L1 & Causal Variable & +$\lambda$=0.01 & threshold > 0.01 \\
L2 & Identity Tracking & +Binding & +8.7\% \\
L3 & Compositional Gen. & +Factorized & error $\sim$ O(N) \\
L4 & Causal Robustness & +Noise Intervention & +24\% \\
L5 & Goal Planning & +Planning Head & Success $\uparrow$ \\
L6 & Scaling Control & +Interaction Density & N$_c$ $\approx$ 4 \\
L7 & Relational Reasoning & +Graph Density & $\lambda \times D > 0.03$ \\
L8 & Physics Adaptation & Match ForceType & Linear $\to$ MLP, Nonlin $\to$ GNN \\
L9 & Theoretical Law & Formula & N $\times$ threshold $>$ critical \\
\end{tabular}
\end{center}

\subsection{Key Finding: Cumulative Effect}

Minimal modifications show cumulative improvement:
\begin{center}
\begin{tabular}{l|c|c}
Stage & MSE & Relative Improvement \\
\hline
Seed & 0.000106 & 0\% \\
+Binding & 0.000097 & +8.7\% \\
+Graph & 0.000091 & +14.9\% \\
\end{tabular}
\end{center}

\section{Constraint Interaction Law}

We discover the fundamental law governing capability emergence:

\begin{equation}
\text{threshold} = \frac{\lambda \times \text{Density}}{\text{Nonlinearity\_exponent}}
\end{equation}

Where:
\begin{itemize}
\item $\lambda$: Binding strength (0-0.5)
\item Density: Graph density (0-1)
\item Nonlinearity\_exponent: 1 for linear (F $\propto$ r), 2 for nonlinear (F $\propto$ r$^2$)
\end{itemize}

Decision rule:
\begin{equation}
\text{OptimalStructure}(N, \text{ForceType}) = 
\begin{cases}
\text{MLP} & \text{if ForceType = linear} \\
\text{GNN} & \text{if ForceType = nonlinear}
\end{cases}
\end{equation}

\section{Physics-Dependent Scaling}

We verify that optimal structure depends on physics nonlinearity:

\begin{center}
\begin{tabular}{c|c|c|c}
Force Type & N=4 & N=8 & N=12 \\
\hline
Linear (F$\propto$r) & MLP +80\% & MLP +85\% & MLP +96\% \\
Nonlinear (F$\propto$r$^2$) & MLP +33\% & MLP +32\% & GNN +24\% \\
\end{tabular}
\end{center}

This confirms: linear systems prefer global MLP, nonlinear systems require graph structure at larger N.

\section{Dimension Exploration}

We systematically explore optimal configurations:

\subsection{Activation Functions}
\begin{center}
\begin{tabular}{l|c}
Activation & MSE \\
\hline
Tanh & 0.000066 \\
ELU & 0.000074 \\
LeakyReLU & 0.000104 \\
ReLU & 0.000109 \\
SiLU & 0.000116 \\
Sigmoid & 0.002142 \\
\end{tabular}
\end{center}

Tanh emerges as optimal for this task.

\subsection{Learning Rates}
\begin{center}
\begin{tabular}{l|c}
Learning Rate & MSE \\
\hline
0.01 & 0.000106 \\
0.001 & 0.000107 \\
0.0001 & 0.001080 \\
0.1 & 0.003453 \\
1e-5 & 0.006170 \\
\end{tabular}
\end{center}

LR=0.01 is optimal.

\subsection{Prediction Horizons}
\begin{center}
\begin{tabular}{l|c}
Horizon & MSE \\
\hline
1-step & 0.001083 \\
2-step & 0.003600 \\
3-step & 0.008315 \\
5-step & 0.021801 \\
10-step & 0.071265 \\
\end{tabular}
\end{center}

Shorter horizons work better.

\section{Cross-Task Transfer}

We test generalization to new tasks:

\begin{center}
\begin{tabular}{l|c|c}
Train & Test & MSE Increase \\
\hline
Force & Force & 1x (baseline) \\
Force & Spring & 55x \\
Spring & Force & 806x \\
Spring & Spring & 1x (baseline) \\
\end{tabular}
\end{center}

Significant transfer gap exists between different physics types.

\section{Limitations}

\begin{enumerate}
\item \textbf{Simple environments}: All experiments in simplified 2D physics simulations
\item \textbf{No real robots}: MuJoCo/real robot validation pending
\item \textbf{Pixel input not tested}: Perception bottleneck not addressed
\item \textbf{Transfer gap}: Cross-task generalization shows significant decay
\end{enumerate}

\section{Conclusion}

We discover that intelligence can emerge from an extremely minimal seed structure through 1-2 parameter modifications. The Constraint Interaction Law provides a theoretical framework for understanding capability emergence. However, generalization to real-world tasks remains an open challenge.

\section{Future Work}

\begin{itemize}
\item MuJoCo/real robot validation
\item Pixel-level perception integration  
\item Multi-modal binding (vision + physics)
\item Longer-horizon prediction
\end{itemize}

\end{document}
